    \documentclass[12pt,a4paper,titlepage,twoside,openright]{article}
    \usepackage{tikz}
    \usetikzlibrary{intersections}
    \usetikzlibrary{shapes,arrows,chains}
    \tikzstyle{line}=[draw] % here
    %%%<
    \usepackage{verbatim}


    \usepackage{amsfonts}
    \usepackage{amsmath}
    \usepackage{cite}
    \usepackage{enumerate}
    \usepackage{flushend}
    \usepackage{graphics,graphicx}
    \usepackage{comment}
    \usepackage{multirow}
    \usepackage{algpseudocode}
    \usepackage{color}
    \usepackage{tikz}
    \usetikzlibrary{shapes,arrows,chains}
    %%%<
    \usepackage{verbatim}

    \DeclareMathOperator*{\argmax}{arg\,max}
    \usepackage{array}
    \renewcommand{\arraystretch}{1.2}
    \usepackage{graphicx}
    \usepackage{subfigure}
    \newtheorem{thm}{Theorem}
    \newtheorem{lmm}{Lemma}
    \newtheorem{cll}{Corollary}
    \newtheorem{dfn}{Definition}
    \newtheorem{prop}{Proposition}
    %\makeatletter

    \begin{document}

    \colorlet{lcfree}{black} 
    \colorlet{lcnorm}{black}
    \colorlet{lccong}{black}

    \providecommand{\cmark}[2][]{\relax}
    \begin{tikzpicture}[%
    >=triangle 60,              % Nice arrows; your taste may be different
    start chain=going below,    % General flow is top-to-bottom
    node distance=10mm and 75mm, % Global setup of box spacing
    every join/.style={norm},   % Default linetype for connecting boxes
    ]


    % ------------------------------------------------- 
    % A few box styles 
    % <on chain> *and* <on grid> reduce the need for manual relative
    % positioning of nodes
    \tikzset{
     base/.style={draw, on chain, on grid, align=center, minimum height=4ex},
     proc/.style={base, rectangle, text width=12em},
     test/.style={base, diamond, aspect=2, text width=8em},
     term/.style={proc, rounded corners},
     % coord node style is used for placing corners of connecting lines
     coord/.style={coordinate, on chain, on grid, node distance=2mm and 12mm},
     % nmark node style is used for coordinate debugging marks
     nmark/.style={draw, cyan, circle, font={\sffamily\bfseries}},
     % -------------------------------------------------
     % Connector line styles for different parts of the diagram
     norm/.style={->, draw, lcnorm},
     free/.style={->, draw, lcfree},
     cong/.style={->, draw, lccong},
     it/.style={font={\small\itshape}}
    }
    % -------------------------------------------------
    % Use join to connect a node to the previous one 
    \node [term, join]  (p0)    {t  =  0\\ Intialize $m_i$, $q_i$, $p_i$,\\ and $q_k$, $p_k$, $a_k$.};
    \node [proc, join] (p1) {Sense AoI and update inforamtion maps.};
    \node [proc, join]      {Detect targets and update inforamtion lists.};
    \node [test, join] (t1) {Is $N_i={\emptyset}$?};
    \node [proc] (p2) {Evaluate the fitness of each target as in Section \ref{sec:Proposed distributed semi-flocking}-A.};
    \node [proc, join] (p3) {Dertermin the next operating mode as in Section \ref{sec:Proposed distributed semi-flocking}-B.};
    \node [test, join] (t2) {Track a target $k$?};
    \node [proc, join=by free]  (p7) {Calculate $f_i^t$ as in Eq. (12).};
    %%%%%%%%%%%%%%%%%%%%%%%%%%%%%%%%%%%%%%%%%%%%%%%%%%%%%%%%%%
    \node [proc,  right=of p1] (p4) {Exchange inforamtion maps with all neighbours.};
    \node [proc, join=by free]   {Transimit information lists to all neighbours.};
    \node [test, join] (t3) {other nodes?};
    \node [proc, join=by free, right=of p2] (p5) {Transimit information lists to these nodes.};
    \node [proc,  right=of t2] (p6) {Calculate $q_{i_s}$ as in Section \ref{sec:Proposed distributed semi-flocking}-C.};
    \node [proc, join=by free, right=of p7]  (p8) {Calculate $f_i^s$ as in Eq. (13).};
    % Some more nodes specifically positioned (we could have avoided this,
    % but try it and you'll see the result is ugly).
    % -------------------------------------------------
    % Now we place the coordinate nodes for the connectors with angles, or
    % with annotations. We also mark them for debugging.
    \node [coord, right=of t1] (c7)  {}; \cmark{7}  



    % -------------------------------------------------
    % A couple of boxes have annotations
    % -------------------------------------------------
    % All the other connections come out of tests and need annotating
    % First, the straight north-south connections. In each case, we first
    % draw a path with a (consistently positioned) annotation node, then
    % we draw the arrow itself.
    \path (t1.south) to node [near start, xshift=1em] {$y$} (p2);
\draw [->,lcnorm] (t1.south) -- (p2);
%\path (t1.east) to node [near start, xshift=1em] {$n$} (p4);
%\draw [->,lcnorm] (t1.east) -- (p4.west);
\draw[->,lcnorm] (t1.east)--++(0:10mm) node[near start, above] {$n$} |-(p4);
\path (t3.south) to node [near start, xshift=1em] {$y$} (p5);
\draw [->,lcnorm] (t3.south) -- (p5);
\path (p5.west) to node [near start, xshift=1em] {} (p2); 
\draw [->,lcnorm] (p5.west) -- (p2);
\path (t2.east) -| node [very near start, yshift=1em] {$n$} (c7); 
\draw [->,lcfree] (t2.east) -- (p6.west);

    % -------------------------------------------------

% -------------------------------------------------
\end{tikzpicture}

\end{document}
